% Рисунок из тетради: канторова лестница — непрерывна, но не абсолютно непрерывна.
\begin{figure}[htbp]
\centering
\begin{tikzpicture}[x=7.2cm, y=3.6cm]
  % оси
  \draw[->] (0,0) -- (1.12,0) node[below right, font=\footnotesize] {$x$};
  \draw[->] (0,0) -- (0,1.15);
  \foreach \y/\l in {0.25/{\tfrac14}, 0.5/{\tfrac12}, 0.75/{\tfrac34}, 1/{1}} {
    \draw[gray!60] (-0.012,\y) -- (0.012,\y);
    \node[anchor=east, font=\scriptsize] at (-0.015,\y) {$\l$};
  }
  \node[anchor=north east, font=\scriptsize] at (0,0) {$0$};
  \node[anchor=north, font=\scriptsize] at (1,0) {$1$};

  % канторова лестница, приближение 3-го уровня
  \draw[very thick, blue!60!black]
    (0,0) -- (0.037037,0.125) -- (0.074074,0.125) -- (0.111111,0.25)
    -- (0.222222,0.25) -- (0.259259,0.375) -- (0.296296,0.375) -- (0.333333,0.5)
    -- (0.666667,0.5) -- (0.703704,0.625) -- (0.740741,0.625) -- (0.777778,0.75)
    -- (0.888889,0.75) -- (0.925926,0.875) -- (0.962963,0.875) -- (1,1);

  % уровни канторова множества
  \begin{scope}[yshift=-1.05cm]
    \foreach \lev/\dy in {1/0, 2/-0.09, 3/-0.18} {
      \node[anchor=east, font=\scriptsize, gray] at (-0.03,\dy) {\lev};
    }
    \node[anchor=east, font=\scriptsize, gray] at (-0.03,0.09) {уровень};
    % уровень 1
    \draw[very thick] (0,0) -- (0.333333,0);
    \draw[very thick] (0.666667,0) -- (1,0);
    % уровень 2
    \draw[very thick] (0,-0.09) -- (0.111111,-0.09);
    \draw[very thick] (0.222222,-0.09) -- (0.333333,-0.09);
    \draw[very thick] (0.666667,-0.09) -- (0.777778,-0.09);
    \draw[very thick] (0.888889,-0.09) -- (1,-0.09);
    % уровень 3
    \foreach \a/\b in {0/0.037037, 0.074074/0.111111, 0.222222/0.259259,
                       0.296296/0.333333, 0.666667/0.703704, 0.740741/0.777778,
                       0.888889/0.925926, 0.962963/1} {
      \draw[very thick] (\a,-0.18) -- (\b,-0.18);
    }
    \node[anchor=west, font=\scriptsize] at (1.03,-0.09)
      {суммарная длина $\left(\tfrac23\right)^{m} \to 0$};
  \end{scope}
\end{tikzpicture}
\caption{Канторова лестница. На канторовом множестве она не задана явно, но
продолжается на весь отрезок по непрерывности. Суммарная длина отрезков $m$-го
уровня равна $(2/3)^m \to 0$, а суммарное приращение функции на них остаётся
равным $1$. Значит, лестница \emph{не} абсолютно непрерывна, хотя непрерывна:
абсолютная непрерывность строго сильнее.}
\label{fig:cantor}
\end{figure}
